%!TEX root = ../main.tex

%%**************************************************************
%% Vorlage fuer Bachelorarbeiten (o.ä.) der Fakultät Technik an der DHBW Stuttgart/Bad Mergentheim
%%
%% Autor: Tobias Dreher, Yves Fischer
%% Datum: 06.07.2011
%%
%% Autor: Michael Gruben
%% Datum: 15.05.2013
%%
%% Autor: Ioanna Theofilou
%% Datum: 10.02.2026
%%**************************************************************

\usepackage{fullpage}

\usepackage[activate]{microtype}

\usepackage[utf8]{inputenc}
\usepackage[T1]{fontenc}

\usepackage[onehalfspacing]{setspace}

\usepackage[
  left=4cm,
  right=2cm,
  top=3cm,
  bottom=3.5cm
]{geometry}

\usepackage{makeidx}

\usepackage[ngerman]{babel}

\usepackage[babel,german=quotes]{csquotes}

% Deckblatt
\usepackage{longtable}
\setlength{\tabcolsep}{10pt}
\renewcommand{\arraystretch}{1.5} 

% Grafiken
\usepackage{graphicx}

% Mathe
%\usepackage{amsmath}
%\usepackage{amssymb}

% Seite drehen
%\usepackage{rotating}
%\usepackage{lscape}

\usepackage{color}
\definecolor{LinkColor}{rgb}{0,0,0.2}
\definecolor{ListingBackground}{rgb}{0.92,0.92,0.92}

% Arbeitstitel (6-12 Worte, praezise, ohne Firmennamen/Abkuerzungen) - mit Betreuer final abstimmen
\newcommand{\pdftitel}{Vergleich verschiedener Lebensdauerversuche von Elektro-Aktuatoren}
\newcommand{\autor}{[Vorname Nachname]}
\newcommand{\arbeit}{Bachelorarbeit}

% Für das Deckblatt + Autor auf den Seiten
\title{\titel}
\author{\autor}
\date{\datum}

% PDF Einstellungen
\usepackage[%
	pdftitle={\pdftitel},
	pdfauthor={\autor},
	pdfsubject={\arbeit},
	pdfcreator={pdflatex, LaTeX with KOMA-Script},
	pdfpagemode=UseOutlines, 
	pdfdisplaydoctitle=true, 
	pdflang=de 
]{hyperref} 

% HTML Links
\hypersetup{
	colorlinks=true,
	linkcolor=LinkColor, 
	citecolor=LinkColor,
	filecolor=LinkColor,
	menucolor=LinkColor,
	urlcolor=LinkColor,
	bookmarksnumbered=true
}

% % Literaturverweise nach Harvard
% \usepackage[dcucite]{harvard}
% \renewcommand{\harvardand}{und}

% Literaturverweise Biblatex
\usepackage[style=authoryear, maxcitenames=1, maxbibnames=99]{biblatex}
\addbibresource{ArbeitBib.bib}

\DefineBibliographyStrings{ngerman}{
  andothers = {\textit{et al.}}
}

% Schriftart
%\usepackage{goudysans}
%\usepackage{lmodern}
%\usepackage{libertine}
\usepackage{palatino} 

\clubpenalty=10000
\widowpenalty=10000
\displaywidowpenalty=10000

% Listings (Quellcode) hier muss für jede Programmiersprache ein eigenes \lstloadlanguages und \lsset gemacht werden. Siehe unten Beispiel.
\usepackage{listings}
\lstloadlanguages{Java}
\lstloadlanguages{python}
\lstset{%
	language=PHP,		 	 % Sprache des Quellcodes
	%numbers=left,           % Zelennummern links
	stepnumber=1,            % Jede Zeile nummerieren.
	numbersep=5pt,           % 5pt Abstand zum Quellcode
	numberstyle=\tiny,       % Zeichengrösse 'tiny' für die Nummern.
	breaklines=true,         % Zeilen umbrechen wenn notwendig.
	breakautoindent=true,    % Nach dem Zeilenumbruch Zeile einrücken.
	postbreak=\space,        % Bei Leerzeichen umbrechen.
	tabsize=2,               % Tabulatorgrösse 2
	basicstyle=\ttfamily\footnotesize, % Nichtproportionale Schrift, klein für den Quellcode
	showspaces=false,        % Leerzeichen nicht anzeigen.
	showstringspaces=false,  % Leerzeichen auch in Strings ('') nicht anzeigen.
	extendedchars=true,      % Alle Zeichen vom Latin1 Zeichensatz anzeigen.
	captionpos=b,            % sets the caption-position to bottom
	backgroundcolor=\color{ListingBackground} % Hintergrundfarbe des Quellcodes setzen.
}

% Glossar
\usepackage[
	nonumberlist, %keine Seitenzahlen anzeigen
	%section,     %im Inhaltsverzeichnis auf section-Ebene erscheinen
	toc,          %Einträge im Inhaltsverzeichnis
]{glossaries}

\usepackage{tocloft}

%Akronyme
\usepackage[printonlyused]{acronym}

% Symbolverzeichnis
\usepackage{nomencl}
\renewcommand{\nomname}{Symbolverzeichnis}
\makenomenclature

% Formelverzeichnis
\usepackage{tocloft}
\newlistof{formulas}{for}{Formelverzeichnis}

% Algorithmusverzeichnis
\usepackage{algorithm}
\usepackage{algpseudocode}
\floatname{algorithm}{Algorithmus}
\renewcommand{\listalgorithmname}{Algorithmusverzeichnis}

% Quellcodeverzeichnis
\usepackage{listings}
\renewcommand\lstlistingname{Quellcode}
\renewcommand\lstlistlistingname{Quellcodeverzeichnis}


% Fussnoten
\usepackage[perpage, hang, multiple, stable]{footmisc}

%Bildpfad
\graphicspath{{images/}}

%nur ein latex-Durchlauf für die Aktualisierung von Verzeichnissen nötig
\usepackage{bookmark}

%Bilder, Tabellen, usw\ldots lassen sich mit H an genau der definierten Stelle platzieren
\usepackage{float}

% für die vertikale Platzierung von Text in Tabellen
\usepackage{array}

% für die Darstellung des Euro-Symbols
\usepackage[right]{eurosym}

% für textumflossene Grafiken
\usepackage{wrapfig}

% eine Kommentarumgebung "k" (Handhabe mit \begin{k}<Kommentartext>\end{k},
% Kommentare werden rot gedruckt). Wird \% vor excludecomment{k} entfernt,
% werden keine Kommentare mehr gedruckt.
\usepackage{comment}
\specialcomment{k}{\begingroup\color{red}}{\endgroup}
%\excludecomment{k}

% individuelle Befehle um auf Kapitel, Sections, Abbildungen und Referenzen allgemein zu verweisen.
\newcommand{\refKap}[1]{Kapitel \ref{#1} auf Seite \pageref{#1} (\nameref{#1})}
\newcommand{\refSec}[1]{Punkt \ref{#1}}
\newcommand{\refAbb}[1]{Abbildung \ref{#1}}
\newcommand{\fullref}[1]{\ref{#1} auf Seite \pageref{#1} (\nameref{#1})}

%moves numbers to the right
\usepackage{fancyhdr}
\pagestyle{fancy}
\fancyhf{}
\fancyfoot[L]{\autor}
\fancyfoot[R]{\thepage} 
\renewcommand{\headrulewidth}{0pt}
\renewcommand{\footrulewidth}{0.6pt}

% Ensure page number appears bottom right even on plain pages
\fancypagestyle{plain}{
  \fancyhf{}
  \fancyfoot[L]{\autor}
  \fancyfoot[R]{\thepage}
  \renewcommand{\headrulewidth}{0pt}
    \renewcommand{\footrulewidth}{0.6pt}
}

% ============= Einstellung für die Größen der verschiedenen Überschriften =============
\setkomafont{chapter}{\LARGE}
\setkomafont{section}{\Large}
\setkomafont{subsection}{\large}
\setkomafont{subsubsection}{\normalsize}
\setkomafont{paragraph}{\normalsize}
\setkomafont{subparagraph}{\small}


% Fragebogen
\usepackage{amssymb}     % Für Kästchen (\square)
\usepackage{enumitem}    % Für flexible Aufzählungen


% PDF einbinden
\usepackage{graphicx} 
\usepackage{subcaption}


% Neues Anhangsverzeichnis definieren
\newlistof{anhangverzeichnis}{apc}{}
\newcommand{\anhang}[1]{
  \refstepcounter{chapter}
  \chapter*{Anhang \thechapter: #1}
  \addcontentsline{apc}{chapter}{Anhang \thechapter: #1}
}